\newpage
\subsection{Performance on more architectures and more challenging tasks}

\begin{table*}[h]
\centering
\caption{{Performance comparison of different compression methods on LLaMA3.1 and Qwen2.5 architectures at 0.2 compression rate across multiple benchmarks. Results are reported for 7 zero-shot tasks (ZeroShot\textsuperscript{7}), long-context understanding (LongEval), code generation (HumanEval), and mathematical reasoning (GSM8K).}}
\label{tab:qwen}
\begin{tabular}{l|l|cccc}
\hline
\textbf{Model} & \textbf{Method} & \textbf{ZeroShot\textsuperscript{7}} & \textbf{LongEval} & \textbf{HumanEval} & \textbf{GSM8K} \\
\hline
\multirow{4}{*}{LLaMA3.1-8B}
 & FP & 0.61 & 0.41 & 0.35 & 0.15 \\
 & ASVD & 0.49 & 0.33 & 0.15 & 0.06 \\
 & SVD-LLM & 0.57 & 0.37 & 0.27 & 0.10 \\
 & \textbf{Ours} & \textbf{0.59} & \textbf{0.39} & \textbf{0.33} & \textbf{0.13} \\
\hline
\multirow{4}{*}{LLaMA3.1-70B}
 & FP & 0.67 & 0.49 & 0.55 & 0.37 \\
 & ASVD & 0.54 & 0.39 & 0.32 & 0.26 \\
 & SVD-LLM & 0.63 & 0.45 & 0.47 & 0.32 \\
 & \textbf{Ours} & \textbf{0.66} & \textbf{0.48} & \textbf{0.52} & \textbf{0.35} \\
\hline
\multirow{4}{*}{Qwen2.5-7B}
 & FP & 0.60 & 0.58 & 0.66 & 0.71 \\
 & ASVD & 0.50 & 0.49 & 0.41 & 0.42 \\
 & SVD-LLM & 0.56 & 0.55 & 0.55 & 0.64 \\
 & \textbf{Ours} & \textbf{0.58} & \textbf{0.57} & \textbf{0.63} & \textbf{0.69} \\
\hline
\multirow{4}{*}{Qwen2.5-32B}
 & FP & 0.65 & 0.79 & 0.55 & 0.72 \\
 & ASVD & 0.52 & 0.61 & 0.37 & 0.44 \\
 & SVD-LLM & 0.61 & 0.73 & 0.46 & 0.63 \\
 & \textbf{Ours} & \textbf{0.64} & \textbf{0.76} & \textbf{0.51} & \textbf{0.68} \\
\hline
\end{tabular}
\end{table*}

Table\ref{tab:qwen} reports the performance of LLaMA3.1 and Qwen2.5 families under a 0.2 compression ratio, comparing three representative low-rank decomposition approaches (ASVD, SVD-LLM, and our proposed SAES-SVD). Across all model scales and tasks, SAES-SVD consistently achieves the best reconstruction fidelity among compressed variants, demonstrating strong robustness in both general-purpose reasoning (ZeroShot$^{7}$), long-context understanding (LongEval), code generation (HumanEval), and mathematical reasoning (GSM8K).

Notably, SAES-SVD narrows the performance gap to the FP baseline much more effectively than existing SVD-based methods. For example, on LLaMA3.1-8B, SAES-SVD retains 96.7\% of FP performance on ZeroShot$^{7}$ and preserves near-lossness on LongEval. Similar trends hold for larger models: on LLaMA3.1-70B and Qwen2.5-32B, SAES-SVD significantly outperforms ASVD and SVD-LLM across all tasks, highlighting its superior scalability to high-capacity architectures.
These results collectively demonstrate that SAES-SVD is a more faithful and stable low-rank approximation method, enabling high compression rates while minimizing degradation across diverse benchmarks.

\subsection{The effect of introducing the mixed-rank configuration}

\begin{table}[h]
\centering
\caption{{Performance of SAES after incorporating the mixed-rank configuration derived from ASVD on the LLaMA-7B model. ZeroShot$^7$ denotes the averaged accuracy over seven zero-shot tasks consistent with the main evaluation protocol.}}
\label{tab:mixrank}
\begin{tabular}{c l c c}
\toprule
\textbf{Compression Ratio} & \textbf{Method} & \textbf{Wiki PPL} & \textbf{ZeroShot$^7$} \\
\midrule
\multirow{3}{*}{0.2}
& ASVD          & 11.14     & 0.43 \\
& SAES-SVD      & 7.17      & 0.50 \\
& SAES+mixrank  & \textbf{6.19} & \textbf{0.51} \\
\midrule
\multirow{3}{*}{0.4}
& ASVD          & 1.00E+03  & 0.30 \\
& SAES-SVD      & 10.42     & 0.41 \\
& SAES+mixrank  & \textbf{7.48} & \textbf{0.47} \\
\bottomrule
\end{tabular}
\end{table}

Table \ref{tab:mixrank} summarizes the effect of introducing the mixed-rank configuration into the SAES framework for the LLaMA-7B model. The mixed-rank strategy is obtained by reusing the adaptive rank distribution discovered by ASVD, while the reconstruction and optimization procedure remains governed by our SAES architecture.
Across both compression ratios(0.2 and 0.4), incorporating mixed rank consistently improves over standard SAES-SVD. At a 0.2 compression ratio, SAES+mixrank achieves the lowest Wiki perplexity(6.19) and the highest zero-shot accuracy (ZeroShot@7 = 0.51), outperforming both ASVD and SAES-SVD. A similar pattern holds at the more aggressive 0.4 compression ratio, where ASVD collapses severely in language modeling, but SAES+mixrank maintains stable performance with a substantial gain in AvgAcc@7(0.47).
These results indicate that adaptive rank allocation and SAES's weight reconstruction are complementary: the mixed-rank configuration provides better structural flexibility, while SAES preserves semantic fidelity during compression. The combination yields a more robust low-rank approximation that scales effectively to higher compression budgets. 


\subsection{Comparison with structured pruning methods}

\begin{table}[t]
\centering
\caption{{Perplexity comparison of SAES-SVD and three representative structured compression methods---LLM-Pruner, SliceGPT, and BlockPruner---on LLaMA-7B under various memory budgets. Lower perplexity indicates better language modeling quality.}}
\label{tab:pruning}
\begin{tabular}{c c c c c}
\toprule
\textbf{Memory} & \textbf{LLM-Pruner} & \textbf{SliceGPT} & \textbf{BlockPruner} & \textbf{SAES-SVD (Ours)} \\
\midrule
10GB & 9.88 & 8.78 & 9.40 & \textbf{7.17} \\
9GB  & 12.21 & 12.73 & 12.76 & \textbf{8.22} \\
8GB  & 18.94 & 16.39 & 19.78 & \textbf{8.96} \\
7GB  & 21.68 & 27.41 & 43.05 & \textbf{10.15} \\
\bottomrule
\end{tabular}
\end{table}

Table \ref{tab:pruning} presents the perplexity of LLaMA-7B compressed using three widely adopted structured pruning and slicing techniques—LLM-Pruner, SliceGPT, and BlockPruner—together with our SAES-SVD method under progressively tighter memory budgets. Across all settings from 10GB down to 7GB, SAES-SVD consistently achieves the lowest perplexity, indicating superior retention of language modeling capability.
When memory reduces to 10GB, SAES-SVD reaches a perplexity of 7.17, surpassing SliceGPT (8.78) and outperforming LLM-Pruner and BlockPruner by a clear margin. As compression becomes more aggressive, the gap further widens: at 8GB and 7GB budgets, structured pruning baselines degrade sharply—BlockPruner even collapsing to a perplexity above 40—while SAES-SVD maintains stable performance with perplexity 8.96 and 10.15, respectively.
This demonstrates that SAES-SVD provides stronger parameter preservation under stringent memory constraints, largely due to its ability to reconstruct informative subspaces without relying on brittle structural sparsity assumptions. As a result, SAES-SVD serves as a more robust alternative to traditional structured compression pipelines, especially in low-memory deployment scenarios.



\subsection{Combination Capability of SAES-SVD and GPTQ}

\begin{table}[h]
\centering
\caption{{Comparison between pure GPTQ-3bit quantization and the equivalent 3bit configuration obtained by combining SAES-SVD with GPTQ on LLaMA2-7B and Qwen2.5-7B. WikiText2 perplexity and ZeroShot@7 represent language modeling quality and the averaged accuracy across seven zero-shot tasks, respectively.}}
\label{tab:gptq}
\begin{tabular}{l l c c}
\toprule
\textbf{Model} & \textbf{Method} & \textbf{Wiki PPL} & \textbf{AvgAcc@7} \\
\midrule
\multirow{2}{*}{\textbf{LLaMA2-7B}}
& GPTQ (3bit)            & 10.17 & 0.39 \\
& SAES+GPTQ (equal 3bit) & \textbf{7.64} & \textbf{0.44} \\
\midrule
\multirow{2}{*}{\textbf{Qwen2.5-7B}}
& GPTQ (3bit)            & 12.29 & 0.51 \\
& SAES+GPTQ (equal 3bit) & \textbf{9.87} & \textbf{0.57} \\
\bottomrule
\end{tabular}
\end{table}


Table \ref{tab:gptq} compares the standard GPTQ 3-bit quantization with our equivalent 3-bit configuration that integrates SAES-SVD into the GPTQ pipeline. The evaluation covers both perplexity on WikiText2 and the averaged accuracy across seven zero-shot tasks for two representative models: LLaMA2-7B and Qwen2.5-7B.
Across both architectures, the combined SAES+GPTQ strategy achieves a substantial improvement over pure GPTQ. For LLaMA2-7B, SAES+GPTQ reduces perplexity from 10.17 to 7.64 and increases AvgAcc@7 from 0.39 to 0.44. A similar trend is observed for Qwen2.5-7B, where perplexity improves from 12.29 to 9.87, accompanied by an accuracy gain from 0.51 to 0.57.
These results demonstrate that SAES-SVD provides a complementary compression effect to GPTQ. By reconstructing a more informative low-rank subspace before quantization, SAES effectively smooths weight distributions and reduces quantization sensitivity, enabling GPTQ to operate at lower bit-widths without incurring substantial performance degradation. Consequently, the SAES+GPTQ pipeline offers a stronger alternative to pure 3-bit quantization while maintaining the same effective model size.




\subsection{The theoretical complexity analysis and comparison of compression time}

We provide a systematic evaluation from two perspectives—(1) \textbf{theoretical complexity analysis} and (2) \textbf{empirical runtime breakdown}—and compare our method (SAES-SVD) with the SVD-LLM baseline.

\subsubsection{Theoretical Complexity Analysis}

\paragraph{Statistics Collection Phase.}
For each layer with input dimension $d_{\mathrm{in}}$, the unfolded calibration set contains $N$ tokens. We maintain only two second-order matrices:
\[
H_\ell = X_\ell X_\ell^\top \in \mathbb{R}^{d_{\mathrm{in}} \times d_{\mathrm{in}}}, 
\quad
\Delta_\ell = (X_\ell^f - X_\ell) X_\ell^\top \in \mathbb{R}^{d_{\mathrm{in}} \times d_{\mathrm{in}}}.
\]
For each mini-batch with $b$ tokens, we perform two matrix multiplications with complexity $\mathcal{O}(d_{\mathrm{in}}^2 b)$, and update the statistics via exponential moving averages. The total complexity over $L$ layers is:
\[
\mathcal{O}\!\left( 2 L d_{\mathrm{in}}^2 N \right),
\]
which matches the order of GPTQ and other second-order methods, and is comparable to the statistics collection cost of SVD-LLM.

\paragraph{Closed-form CEALC Solution for a Single Layer.}
For a linear layer $W_\ell \in \mathbb{R}^{d_{\mathrm{out}} \times d_{\mathrm{in}}}$ with target rank $r_\ell$, the main computational costs of CEALC are:

\begin{itemize}
    \item \textbf{Whitening matrix construction:} Compute $L_\ell = (H_\ell + \lambda I)^{-1/2}$ using a Cholesky-based procedure, with cost $\mathcal{O}(d_{\mathrm{in}}^3)$. This cost is shared by all second-order compression methods.
    \item \textbf{Target matrix construction:} Build $G_\ell = W_\ell(H_\ell + \beta_\ell \Delta_\ell)L_\ell$, requiring several matrix multiplications with total complexity $\mathcal{O}(d_{\mathrm{out}} d_{\mathrm{in}}^2)$.
    \item \textbf{Truncated SVD:} Perform a rank-$r_\ell$ truncated SVD on $G_\ell$, with complexity $\mathcal{O}(d_{\mathrm{out}} d_{\mathrm{in}} r_\ell)$.
\end{itemize}

\paragraph{ACES Adaptive Coefficient Selection.}
During compression, ACES introduces additional computations:

\begin{itemize}
    \item \textbf{Shared whitened matrices:} Construct
    \[
    S_\ell = W_\ell H_\ell L_\ell, \quad D_\ell = W_\ell \Delta_\ell L_\ell,
    \]
    each requiring $\mathcal{O}(d_{\mathrm{out}} d_{\mathrm{in}}^2)$.
    \item \textbf{Subspace extraction:} Perform a rank-$r_\ell$ truncated SVD on $S_\ell$ with complexity $\mathcal{O}(d_{\mathrm{out}} d_{\mathrm{in}} r_\ell)$.
    \item \textbf{Closed-form $\beta_\ell$ selection:} Using Theorem~4.2, compute a small number of Frobenius norms and inner products $(a,b,c,A,B,C)$ within the fixed subspace and solve a univariate quadratic equation in closed form. This step is negligible relative to the matrix operations.
\end{itemize}

Thus, compared with CEALC alone, ACES introduces only one extra truncated SVD per layer and inexpensive scalar operations.

\paragraph{Fine-tuning Phase.}
SVD-LLM requires approximately 3.5 hours of fine-tuning for LLaMA-7B, and even more for larger models. In contrast, \textbf{SAES-SVD requires no fine-tuning} and operates purely as a post-hoc compression method.

\paragraph{Summary.}
SAES-SVD and SVD-LLM share the same order of complexity in statistics collection and SVD-based factorization. However, by eliminating the fine-tuning phase, SAES-SVD substantially reduces the overall compression time.

\subsubsection{Empirical Runtime Breakdown and Comparison with the Baseline}

\begin{table}[h]
\centering
\caption{{Comparison of SAES-SVD and SVD-LLM regarding statistics time, compression time, and fine-tuning time on LLaMA-7B.}}
\begin{tabular}{lccccc}
\toprule
\textbf{Method} & \textbf{Data Collection} & \textbf{Compression} & \textbf{Fine-tuning} & \textbf{Total} & \textbf{Reduction} \\
\midrule
SVD-LLM  & 0.45h & 0.15h & 3.5h & 4.1h & -- \\
SAES-SVD & 0.30h & 0.20h & 0    & \textbf{0.5h} & \textbf{-88\%} \\
\bottomrule
\end{tabular}
\end{table}

We further provide in Table~R3-1 an empirical breakdown of the compression time for SAES-SVD and SVD-LLM on LLaMA-7B, measured on an NVIDIA A6000 GPU. The results show:

\begin{itemize}
    \item \textbf{More efficient statistics collection.}  
    Although the theoretical complexity is similar, SAES-SVD uses a smaller calibration set (128 vs.\ 256 samples), reducing statistics collection time to about 67\% of that of SVD-LLM.
    \item \textbf{Controlled compression overhead.}  
    The SVD compression phase of SAES-SVD is slightly more expensive due to ACES, but remains within the same complexity order.
    \item \textbf{Elimination of fine-tuning.}  
    SAES-SVD is a post-training method and does not rely on any fine-tuning, whereas SVD-LLM requires more than 3.5 hours of fine-tuning.
\end{itemize}

Overall, on LLaMA-7B, SAES-SVD achieves a total compression wall-clock time less than one-eighth of that of SVD-LLM, while obtaining better accuracy retention across multiple tasks. This combination of higher accuracy and lower end-to-end compression time highlights the practical efficiency of SAES-SVD for real-world deployment.


\color{black}